% Actividades bibliográficas habilitadas por \ifHAFOPedagogy. Se invocan en el punto
% matemático correspondiente; el archivo no produce contenido por sí solo.

\newcommand{\HAFOProblemStateMemory}{\ifHAFOPedagogy
\begin{bloqueestudio}{Ejercicios}
Para una matriz \(A\in\R^{n\times n}\), resuelva primero
\(\dot x=Ax\) por la serie de la exponencial. Multiplique
\(e^{tA}e^{sA}\) y agrupe los términos de igual potencia de \(A\) para
demostrar \(e^{(t+s)A}=e^{tA}e^{sA}\). Después sustituya la serie
\(E_q(At^q)\) en \(x=x_0+I^qAx\), integrando cada monomio con Beta.
En el caso escalar, compare los términos de orden \(h^q\) en
\(E_q(\lambda(2h)^q)\) y \(E_q(\lambda h^q)^2\): la diferencia
\(\lambda(2^q-2)h^q/\Gamma(q+1)\) identifica por qué falla la
composición puntual cuando \(0<q<1\). Reconstruya el término libre
heredado al separar Volterra en \([0,t]\) y \([t,t+s]\).
\end{bloqueestudio}
\fi}

\newcommand{\HAFOProblemStabilityBasins}{\ifHAFOPedagogy
\begin{bloqueestudio}{Ejercicios}
Elija un Chua entero y uno fraccionario. Resuelva \(F(e)=0\), calcule
\(DF(e)\) y localice su espectro. En orden entero compare la parte real
de los autovalores con cero; en Caputo compare sus argumentos con
\(q\pi/2\). Para un equilibrio estable, formule la vecindad de datos que
converge a él y relacione esa vecindad con la cuenca de otro destino.
Construya una situación en la que un equilibrio inestable tenga una
variedad inestable que alcance un atractor autoexcitado. Esta construcción
muestra por qué la inestabilidad de todos los equilibrios no implica
ocultedad: la condición relevante es la intersección de las cuencas con sus
vecindades, no sólo el espectro.
\end{bloqueestudio}
\fi}

\newcommand{\HAFOProblemReturn}{\ifHAFOPedagogy
\begin{bloqueestudio}{Ejercicios}
Para PLL o MAVPD, elija \(\Sigma=\{h=0\}\) y calcule
\(Dh\,f\) en los cruces. Escriba \(P(z)=\varphi_{\tau(z)}(z)\), derive
la identidad \(h(P(z))=0\) y despeje \(D\tau\). Sustituya ese resultado
en \(DP=D\varphi+f\,D\tau\) para obtener la proyección transversal.
Sobre un abierto acotado de una carta, estudie el grado de \(P-I\) y la
exclusión de ceros en su frontera. Compare existencia del punto fijo con
estabilidad según el espectro de \(DP\). Finalmente, levante la sección a
perfiles de Caputo y escriba \(\mathcal P(m)=S_{\tau(m)}m\); identifique
qué información histórica se pierde al conservar sólo su valor en cero.
\end{bloqueestudio}
\fi}

\newcommand{\HAFOProblemChaos}{\ifHAFOPedagogy
\begin{bloqueestudio}{Ejercicios}
Compare dos trayectorias mediante los mismos observables: una de orden
entero y otra de Caputo. Defina acotación, recurrencia y atracción con sus
límites temporales, y escriba por separado una condición de caos sobre un
conjunto invariante y la condición de ocultedad de su cuenca. El sistema
lineal estable proporciona una trayectoria acotada sin caos; un ciclo
repulsor proporciona recurrencia sin atracción. Use estos ejemplos para
examinar las implicaciones entre las propiedades. En Caputo, compare
además recurrencia del estado y del perfil: una FFT o un crecimiento de
perturbaciones en una ventana describe observables distintos de la cuenca.
\end{bloqueestudio}
\fi}

\newcommand{\HAFOProblemLocalization}{\ifHAFOPedagogy
\begin{bloqueestudio}{Ejercicios}
En un mismo campo, produzca dos semillas mediante función descriptiva,
continuación, puntos perpetuos o bisección de borde. Escriba la ecuación
algebraica que define cada semilla y calcule su residual. Integre ambas con
el mismo operador y terminal inferior y compare sus conjuntos de retornos.
Una semilla es un dato inicial; su trayectoria depende del problema
integral, y el destino asintótico se define mediante un conjunto límite.
Estudie la atracción de datos próximos y la intersección con vecindades de
equilibrios para relacionar ese destino con un atractor oculto. En Caputo,
distinga las semillas insertadas como perfiles constantes de los estados
alcanzados mediante una historia heredada.
\end{bloqueestudio}
\fi}

\newcommand{\HAFOThesisDefense}{\ifHAFOPedagogy
\begin{bloqueestudio}{Ejercicios}
Considere un campo de Lur'e \(f(x)=Ax+b\psi(c^{\mathsf T}x)\).
Resuelva \(f(e)=0\), derive
\(Df(e)=A+b\psi'(c^{\mathsf T}e)c^{\mathsf T}\) y compare el criterio
espectral entero con Matignon. A continuación, aplique la transformada de
Laplace al sistema lineal con entrada y salida \(c^{\mathsf T}x\) para
deducir \(W_q(s)=c^{\mathsf T}(s^qI-A)^{-1}b\). Relacione el cierre de
balance armónico con una semilla y, mediante un retorno, con una órbita.
Para el caso fraccionario, separe la integral de Volterra durante una rampa
del campo y determine qué perfil hereda la dinámica final. Compare la
condición \(Df(p)f(p)=0\) con un evento FPP--H y termine formulando la
atracción y la separación de cuenca del destino. Esta cadena une las
construcciones algebraicas, la evolución dinámica y la geometría global.
\end{bloqueestudio}
\fi}

\newcommand{\HAFOCardTopology}{\ifHAFOPedagogy
\begin{bloqueestudio}{Lecturas}
\begin{itemize}[nosep]
  \item \emph{Topología y geometría diferencial con aplicaciones a la física}
  \cite{NahmadAchar2022Topologia}: coordenadas adaptadas y curvas integrales,
  §§11.3--11.5, pp.~124--125; derivadas de Lie, §§11.6--11.14,
  pp.~125--130.
\end{itemize}
\end{bloqueestudio}
\begin{bloqueestudio}{Ejercicios}
Resuelva §§11.15--11.18, pp.~131--133. Calcule
\(L_fV=\nabla V\cdot f\) para \(V(x)=\|x\|^2/2\).

Demuestre que \(\mathbb R\times S^1\) es el espacio
de fases del PLL y defina una distancia con
\(\theta\sim\theta+2\pi\).
\end{bloqueestudio}
\begin{bloqueestudio}{Lecturas complementarias}
Mismo libro: conexiones lineales, pp.~135--142.
\end{bloqueestudio}
\fi}

\newcommand{\HAFOCardFlow}{\ifHAFOPedagogy
\begin{bloqueestudio}{Lecturas}
\begin{itemize}[nosep]
  \item \emph{Teoría geométrica de ecuaciones diferenciales I}
  \cite{JaurezEtAl2021TGEDI}: cap.~1, p.~1; cap.~3, pp.~53--80; cap.~4,
  pp.~93--155.
  \item \emph{Teoría geométrica de ecuaciones diferenciales II}
  \cite{JaurezEtAl2021TGEDII}: cap.~8, pp.~69--98; existencia y unicidad,
  p.~73; primera variación, p.~74; dependencia del dato inicial, p.~92.
  \item Strogatz, \emph{Nonlinear Dynamics and Chaos}
  \cite{Strogatz2015Nonlinear}: PDF~30--50 y 140--212.
\end{itemize}
\end{bloqueestudio}
\begin{bloqueestudio}{Ejercicios}
Resuelva Strogatz 2.4.1--2.4.8, 2.5.1--2.5.5 y 2.6.2,
PDF~57--59. Para \(\dot x=Ax\), demuestre
\(\varphi_t(x_0)=e^{tA}x_0\) y
\(\varphi_{t+s}=\varphi_t\circ\varphi_s\).

Calcule el jacobiano y el espectro del origen en un
caso entero. Explique qué falla en
Hartman--Grobman si hay valores propios sobre el eje imaginario.
\end{bloqueestudio}
\begin{bloqueestudio}{Lecturas complementarias}
\emph{Métodos topológicos en el estudio de las ecuaciones diferenciales no
lineales} \cite{Amster2021MetodosTopologicos}: apéndice A, pp.~319--327.
\end{bloqueestudio}
\fi}

\newcommand{\HAFOCardBasins}{\ifHAFOPedagogy
\begin{bloqueestudio}{Lecturas}
\begin{itemize}[nosep]
  \item \emph{Exploring Chaos} \cite{Davies2004ExploringChaos}:
  sistemas bidimensionales y fronteras de cuenca, pp.~107--148; §4.7,
  p.~134.
  \item Alligood, Sauer y Yorke, \emph{Chaos}
  \cite{AlligoodSauerYorke1996Chaos}: cuencas, PDF~146--165; fronteras,
  PDF~181--209.
\end{itemize}
\end{bloqueestudio}
\begin{bloqueestudio}{Ejercicios}
Resuelva Alligood--Sauer--Yorke, \emph{Challenge} 4,
PDF~199--201.

Para dos atractores \(A_1,A_2\), escriba
\(\mathcal B(A_i)\), su frontera común y un par de puntos para una bisección
de borde. Indique qué resultado sería sólo local y finito.
\end{bloqueestudio}
\begin{bloqueestudio}{Lecturas complementarias}
Alligood--Sauer--Yorke: \emph{Challenge} 10, PDF~444--451.
\end{bloqueestudio}
\fi}

\newcommand{\HAFOCardDegree}{\ifHAFOPedagogy
\begin{bloqueestudio}{Lecturas}
\emph{Métodos topológicos} \cite{Amster2021MetodosTopologicos}: índice de
curvas, pp.~74--82; disparo, p.~84; Brouwer y Poincaré, p.~99; grado de
Brouwer, pp.~103--145.
\end{bloqueestudio}
\begin{bloqueestudio}{Ejercicios}
Resuelva los problemas del cap.~3, pp.~153--155, y calcule
\(\deg(I,B,0)\), \(\deg(-I,B,0)\).

Para un mapa de retorno \(P\),
formule \(F(z)=P(z)-z\) y enumere las hipótesis necesarias para usar el grado.
\end{bloqueestudio}
\fi}

\newcommand{\HAFOCardLinearGeometry}{\ifHAFOPedagogy
\begin{bloqueestudio}{Lecturas}
\emph{Teoría geométrica de ecuaciones diferenciales II}
\cite{JaurezEtAl2021TGEDII}: primera variación, p.~74; dependencia del dato
inicial, p.~92; Grobman--Hartman, p.~94.
\end{bloqueestudio}
\begin{bloqueestudio}{Ejercicios}
Resuelva Strogatz 6.5.11, PDF~206, y 6.8.6--6.8.7,
PDF~213--214. Derive \(\dot\Phi=Df(x(t))\Phi\), \(\Phi(0)=I\).

Para un equilibrio de Chua, MAVPD, Kalman--Fitts o
PLL, separe subespacios estable, inestable y central.
\end{bloqueestudio}
\fi}

\newcommand{\HAFOCardPoincare}{\ifHAFOPedagogy
\begin{bloqueestudio}{Lecturas}
\begin{itemize}[nosep]
  \item \emph{Teoría geométrica de ecuaciones diferenciales II}
  \cite{JaurezEtAl2021TGEDII}: rectificación y Poincaré, pp.~113--115;
  Floquet y monodromía, pp.~121--141; Poincaré--Bendixson, p.~150.
  \item Strogatz \cite{Strogatz2015Nonlinear}: PDF~215--260 y 261--325.
\end{itemize}
\end{bloqueestudio}
\begin{bloqueestudio}{Ejercicios}
Resuelva Strogatz 7.3.1--7.3.3, PDF~250--251, y demuestre
que una órbita periódica da un punto fijo del retorno.

Construya una sección transversal para PLL o MAVPD y
explique un multiplicador de módulo mayor que uno.
\end{bloqueestudio}
\begin{bloqueestudio}{Lecturas complementarias}
Strogatz 8.4.3--8.4.4, PDF~313--314.
\end{bloqueestudio}
\fi}

\newcommand{\HAFOCardChaos}{\ifHAFOPedagogy
\begin{bloqueestudio}{Lecturas}
\begin{itemize}[nosep]
  \item Strogatz \cite{Strogatz2015Nonlinear}: Lorenz, PDF~326--371;
  Lyapunov y mapas, PDF~372--421; atractores extraños, PDF~446--476.
  \item Alligood--Sauer--Yorke \cite{AlligoodSauerYorke1996Chaos}:
  Lyapunov y herradura, PDF~211--247; atractores, PDF~248--289; caos en
  flujos, PDF~376--415; variedades y crisis, PDF~416--463.
\end{itemize}
\end{bloqueestudio}
\begin{bloqueestudio}{Ejercicios}
Resuelva Strogatz 9.2.2--9.2.3 y 9.3.8--9.3.9,
PDF~366--368; 10.3.7--10.3.8 y 10.5.1--10.5.4, PDF~413--417.

Defina por separado acotación, recurrencia, atracción,
caos y ocultedad. Construya un contraejemplo para cada implicación inválida.

Repita la tabla de implicaciones para una trayectoria de
Caputo: sustituya periodicidad exacta por recurrencia proyectada y declare qué
información depende de la historia.
\end{bloqueestudio}
\begin{bloqueestudio}{Lecturas complementarias}
\begin{itemize}[nosep]
  \item \emph{Exploring Chaos} \cite{Davies2004ExploringChaos}, pp.~44--60,
  65--106 y 184--193.
  \item \emph{Chaos and Time-Series Analysis}
  \cite{Sprott2003ChaosTimeSeries}, pp.~72--183, 213--272 y 302--328.
  \item Ott \cite{Ott2002Chaos}: conjuntos caóticos no atractores y
  fronteras, pp.~168--210; crisis y transitorios, pp.~304--344.
  \item T{\'e}l--Gruiz \cite{TelGruiz2006Chaotic}: sillas, escapes,
  fronteras y caos transitorio, pp.~191--225.
  \item Devaney \cite{Devaney1989Chaotic}: PDF aprox.~29--33, 70,
  110--129.
  \item Hilborn \cite{Hilborn2000Chaos}: PDF aprox.~38--80 y 162--220.
\end{itemize}
\end{bloqueestudio}
\fi}

\newcommand{\HAFOCardFractal}{\ifHAFOPedagogy
\begin{bloqueestudio}{Lecturas complementarias}
\begin{itemize}[nosep]
  \item Falconer \cite{Falconer2014Fractal}: dimensión de caja,
  PDF~59--75; Hausdorff, PDF~76--97; sistemas dinámicos, PDF~238--264.
  \item Barnsley, \emph{Fractals Everywhere}: métricas y topología,
  PDF~12--48; contracciones, PDF~49--129; dimensión, PDF~186--219.
  \item Mandelbrot: ejemplos geométricos, PDF~33 y 277; Falconer desarrolla sus fundamentos de dimensión.
\end{itemize}
\end{bloqueestudio}
\begin{bloqueestudio}{Ejercicios}
Resuelva un ejercicio de dimensión de caja de Falconer,
cap.~2; contraste después con el solucionario, PDF~1--21.

Calcule la dimensión de una frontera de cuenca en dos
refinamientos y explique por qué el resultado no demuestra Wada.
\end{bloqueestudio}
\fi}

\newcommand{\HAFOCardFractional}{\ifHAFOPedagogy
\begin{bloqueestudio}{Lecturas}
\begin{itemize}[nosep]
  \item Oliveira \cite{Oliveira2019SolvedFractional}: Gamma, PDF~30--41;
  Mittag--Leffler, PDF~81--124; Laplace, PDF~130--137; GL, RL y Caputo,
  PDF~179--189.
  \item Herrmann \cite{Herrmann2014Fractional}: PDF~28--53 y 64--89.
  \item Das \cite{Das2020Kindergarten}: construcción y relación RL--Caputo,
  PDF~156--163; evaluación discreta GL/RL, PDF~183--185.
  \item \emph{The Analysis of Fractional Differential Equations}
  \cite{Diethelm2010AnalysisFDE}: pp.~13--73.
\end{itemize}
\end{bloqueestudio}
\begin{bloqueestudio}{Ejercicios}
Oliveira: ejercicios 13--16 y 30--32,
PDF~142--172; ejercicios RL--Caputo 12, 14, 15, 20, 25, 32, 33 y 40,
PDF~191--224. Obtenga
\(\mathcal L\{t^{\beta-1}E_{\alpha,\beta}(\lambda t^\alpha)\}
=s^{\alpha-\beta}/(s^\alpha-\lambda)\).

Derive la transformada de Caputo con sus datos
iniciales y compare RL, GL y Caputo en \(t^\mu\).

Resuelva \(\dot x=\lambda x\) y
\({}^{\mathrm C}D^q x=\lambda x\). Compare \(e^{\lambda t}\) con
\(E_q(\lambda t^q)\), la ley de semigrupo y los criterios espectrales entero y
de Matignon. Después aproxime \({}^{\mathrm C}D^qx=-ax\) con GL/RL y resuelva
el PVI con ABM--PECE; explique por qué discretizar un operador no equivale
automáticamente a resolver el mismo PVI Caputo.
\end{bloqueestudio}
\begin{bloqueestudio}{Lecturas complementarias}
Oldham--Spanier \cite{OldhamSpanier1974Fractional}: p.~16,
pp.~47--60, 69--94 y 136. Herrmann: ejercicios 3.3, 5.1--5.2 y 6.1;
soluciones PDF~423, 428--429 y 432.
\end{bloqueestudio}
\fi}

\newcommand{\HAFOCardMemory}{\ifHAFOPedagogy
\begin{bloqueestudio}{Lecturas}
\begin{itemize}[nosep]
  \item Herrmann \cite{Herrmann2014Fractional}: no localidad y memoria,
  PDF~108--126.
  \item Das \cite{Das2020Kindergarten}: contexto analítico, PDF~41--47;
  ramas y cortes, PDF~43--44 y apéndice detallado PDF~503--507;
  inicialización, kernels y memoria desvaneciente RL, PDF~109--110 y 119--135;
  historia e inicialización RL, PDF~247--255; GL y composición, PDF~256--259;
  relación RL--Caputo, PDF~259--264; inicialización Caputo, PDF~264--273;
  seno con terminal inferior, PDF~273--274;
  Laplace con inicialización, PDF~275--291; resolución de EFD, PDF~292--310;
  EFD lineales e integrales, PDF~333--357; ecuaciones secuenciales, matrices y
  funciones alfa-exponenciales, PDF~358--371; sistemas multivariables,
  PDF~372--375.
  \item Oliveira \cite{Oliveira2019SolvedFractional}: ejercicio 26 sobre
  efecto de memoria, PDF~272--273.
  \item Diethelm \cite{Diethelm2010AnalysisFDE}: PVI Caputo de término único,
  existencia, unicidad, perturbaciones, regularidad y estabilidad,
  pp.~85--126 y 157--162.
\end{itemize}
\end{bloqueestudio}
\begin{bloqueestudio}{Ejercicios}
Herrmann: ejercicio 8.1, PDF~126; solución, PDF~436.
Derive la ecuación integral de Volterra.

Construya dos historias con el mismo valor terminal y
pasados distintos; muestre por qué reiniciar cambia el PVI y por qué un retorno
puntual no es, en general, un mapa de estado. Continúe una trayectoria con
historia heredada y reinicie otra desde el mismo punto; identifique el término
libre de Volterra que cambia.

Para \(f(x)=-x\), verifique que
\(m_c(\theta)=c[1+\theta^q/\Gamma(q+1)]\) es fijo bajo el semiflujo de
Volterra. Use \(c\to0\) para explicar por qué la atracción de reinicios no
implica atracción de una vecindad abierta de todos los términos libres.
Distinga la familia local \(\mathcal D_A\), el universo global
\(\mathcal D\) y un espacio invariante \(X_{\mathrm{adm}}\) cuya
construcción debe justificarse.
\end{bloqueestudio}
\begin{bloqueestudio}{Lecturas complementarias}
Matignon \cite{Matignon1996} para el criterio angular.
Li--Chen--Podlubny \cite{LiChenPodlubny2010} para estabilidad directa y
Mittag--Leffler; no es una prueba de caos ni ocultedad.
\end{bloqueestudio}
\fi}

\newcommand{\HAFOCardExistence}{\ifHAFOPedagogy
\begin{bloqueestudio}{Lecturas}
\emph{Métodos topológicos} \cite{Amster2021MetodosTopologicos}: función de
Green, p.~229; Schauder, pp.~245--267; Leray--Schauder, pp.~273--300.
\end{bloqueestudio}
\begin{bloqueestudio}{Ejercicios}
Resuelva Amster 7.9.1--7.9.2, pp.~313--316.

Escriba el PVI de Caputo como
\(x=x_0+I^qf(x)\). Identifique continuidad,
compacidad y cota a priori requeridas por Schauder o Leray--Schauder; declare
el abierto acotado y la exclusión de puntos fijos en su frontera. Distinga
existencia de atracción y ocultedad.
\end{bloqueestudio}
\fi}

\newcommand{\HAFOCardLure}{\ifHAFOPedagogy
\begin{bloqueestudio}{Lecturas}
Slotine--Li \cite{SlotineLi1991}: realimentación y estabilidad,
PDF~119--175; función descriptiva, PDF~176--209. Para atractores ocultos:
Leonov--Kuznetsov \cite{LeonovKuznetsovVagaitsev2011,LeonovKuznetsov2013}
y localización de Chua \cite{KuznetsovEtAl2017}.
\end{bloqueestudio}
\begin{bloqueestudio}{Ejercicios}
Resuelva Slotine--Li, cap.~5, PDF~207--209, y obtenga una
función descriptiva desde su integral de Fourier.

Derive
\(W_q(s)=c^{\mathsf T}(s^qI-A)^{-1}b\). Para ello escriba
\((s^qI-A)\widehat x=b\widehat u\) con datos iniciales nulos y
proyecte la solución mediante \(c^{\mathsf T}\). Fije la rama de
\(s^q\) y compruebe que \(W_-=-W_+\) transforma
\(1-W_+N=0\) en \(1+W_-N=0\).
\end{bloqueestudio}
\fi}

\newcommand{\HAFOCardMachado}{\ifHAFOPedagogy
\begin{bloqueestudio}{Lecturas}
Machado \cite{Machado2015FDF} y Barbosa--Machado
\cite{BarbosaMachado2002BacklashImpact}.
\end{bloqueestudio}
\begin{bloqueestudio}{Ejercicios}
Para \(N=re^{i\theta}\), \(r>0\), use
\(\log N=\log r+i(\theta+2\pi k)\) para construir
\(N^\mu=r^\mu e^{i\mu(\theta+2\pi k)}\). Calcule la rama principal
en \(\mu=1/2\) y compare las fases de dos elecciones de \(k\).

Para \(N_\mu=N^\mu\), calcule por separado el
residuo auxiliar y el balance físico; determine cuándo la potencia sólo
diversifica semillas.

Obtenga dos semillas de ramas distintas, intégrelas bajo
una misma especificación numérica, entera o Caputo, y compare destino y contacto con las cuencas
conocidas.
\end{bloqueestudio}
\fi}

\newcommand{\HAFOCardContinuation}{\ifHAFOPedagogy
\begin{bloqueestudio}{Lecturas}
Diethelm \cite{Diethelm2010AnalysisFDE}: ABM, pp.~195--210; esquema numérico
multi-término, pp.~211--224, con ejercicio en p.~225. Comparación de
aproximación frecuencial, predictor--corrector y Adomian en
Sun--He--Wang \cite{SunHeWang2022FractionalChaos}, pp.~27--74.
Predictor--corrector original \cite{DiethelmFordFreed2002}.
\end{bloqueestudio}
\begin{bloqueestudio}{Ejercicios}
Derive los pesos predictor y corrector desde Volterra y
aplíquelos a \({}^{\mathrm C}D^qx=-x\).

Escriba una homotopía \(f_\lambda\), transporte una
semilla y demuestre por qué reiniciar la memoria en cada \(\lambda\) define
otro problema. Compare dos algoritmos únicamente después de igualar operador,
terminal inferior, historia, paso y tolerancia.
\end{bloqueestudio}
\fi}

\newcommand{\HAFOCardPerpetual}{\ifHAFOPedagogy
\begin{bloqueestudio}{Lecturas}
Prasad \cite{Prasad2015Perpetual}; Dudkowski--Prasad--Kapitaniak
\cite{DudkowskiPrasadKapitaniak2015,DudkowskiPrasadKapitaniak2017,
DudkowskiPrasadKapitaniak2018};
contraejemplos \cite{NazarimehrSaediJafari2017,NagChowdhuryGhosh2020};
comparación de métodos \cite{GuanXie2025}. Das, Oliveira y Herrmann sustentan
operadores, inicialización y memoria, pero no formulan puntos perpetuos
fraccionarios: FPP--A y FPP--H se presentan como formulaciones metodológicas,
no como resultados de esas fuentes.
\end{bloqueestudio}
\begin{bloqueestudio}{Ejercicios}
Para \(\dot x=f(x)\), derive
\(\ddot x=Df(x)f(x)\) y resuelva los PP de dos campos polinomiales sencillos.

Resuelva \(Df(p)f(p)=0\), \(f(p)\ne0\), en uno de los
sistemas analizados. Desde Volterra obtenga
\(h^{2q}Df(p)f(p)/\Gamma(2q+1)\). Parta de
\(x(h)-p=h^qf(p)/\Gamma(q+1)+O(h^{2q})\), sustituya
\(f(x(s))=f(p)+Df(p)(x(s)-p)+O(\|x(s)-p\|^2)\) en Volterra e
integre \((h-s)^{q-1}s^q\) con la función Beta. Este cálculo requiere
regularidad suficiente del campo para el resto y no usa una regla de cadena
fraccionaria. Explique por qué cancelar el coeficiente de arranque no
define una aceleración global de la trayectoria.

Para \(f(x)=Ax+d\), compare FPP--A y FPP--H. Después,
en un campo no lineal, evalúe el evento FPP--H con toda la historia y contraste
continuación heredada contra reinicio. Repita L1 con \(h,h/2,h/4\) y terminal
inferior común; un mínimo observado gráficamente no constituye un evento
validado.
\end{bloqueestudio}
\fi}

\newcommand{\HAFOCardGeometric}{\ifHAFOPedagogy
\begin{bloqueestudio}{Lecturas}
Curvas conectantes \cite{GilmoreEtAl2010Connecting,GuanXie2021Connecting};
superficies críticas \cite{GodaraEtAl2018CriticalSurfaces}; cuencas y Wada
\cite{DazaEtAl2016BasinEntropy,DazaEtAl2015Wada}.
\end{bloqueestudio}
\begin{bloqueestudio}{Ejercicios}
Para un campo plano, calcule \(Df\), \(Df\,f\),
\(\det Df\) y sus conjuntos de nivel cero.

Repita el cálculo en PLL o MAVPD. Formule una
bisección entre dos destinos con métrica escalada; añada un tercer destino y
describa el fallo.
\end{bloqueestudio}
\fi}

\newcommand{\HAFOCardProtocols}{\ifHAFOPedagogy
\begin{bloqueestudio}{Lecturas}
Datseris--Wagemakers \cite{datseris2022} para estimación de cuencas; entropía de
cuenca \cite{DazaEtAl2016BasinEntropy}; clasificación de cuencas
\cite{SprottXiong2015Basins}; Wada \cite{DazaEtAl2015Wada}.
\end{bloqueestudio}
\begin{bloqueestudio}{Ejercicios}
Para una órbita conocida, calcule retorno, período dominante
y variación al reducir el paso.

Construya una matriz
semilla\(\times\)orden--memoria\(\times\)paso\(\times\)horizonte
\(\times\)solucionador. Separe
destino, atracción, caos, contacto con equilibrios y ocultedad.
\end{bloqueestudio}
\fi}

\newcommand{\HAFOCardReliability}{\ifHAFOPedagogy
\begin{bloqueestudio}{Lecturas}
\begin{itemize}[nosep]
  \item Vallejo--Sanju{\'a}n \cite{VallejoSanjuan2017Predictability}:
  error y \emph{shadowing}, pp.~13--22; FTLE y regímenes, pp.~25--87;
  predictibilidad, pp.~91--126; algoritmo, pp.~129--136.
  \item Datseris--Parlitz \cite{DatserisParlitz2022Nonlinear}: Lyapunov y
  horizonte de predictibilidad, pp.~37--48; bifurcaciones y dimensión,
  pp.~53--85; algoritmo de Lyapunov, pp.~211--216.
  \item Ott \cite{Ott2002Chaos}, pp.~168--210 y 304--344, y
  T{\'e}l--Gruiz \cite{TelGruiz2006Chaotic}, pp.~191--225, para sillas,
  escapes, fronteras y crisis.
  \item Caso primario de transitorio oculto
  \cite{Danca2016HiddenTransient}; exponentes fraccionarios y su algoritmo
  \cite{DancaKuznetsov2018Lyapunov}; memoria e inicialización
  \cite{Sabatier2021InfiniteMemory}.
\end{itemize}
\end{bloqueestudio}
\begin{bloqueestudio}{Ejercicios}
Para una aproximación \(\widetilde x\), escriba
\(r=\widetilde x'-f(\widetilde x)\). Reste su ecuación de la solución
exacta y derive
\(\|e(t)\|\le e^{Lt}\|e(0)\|+\int_0^te^{L(t-s)}\|r(s)\|\,\dd s\).
Compare esta cota Lipschitz con la estimación local
\(T\approx\lambda^{-1}\log(\Delta/\delta)\) cuando un error inicial
\(\delta\) crece como \(e^{\lambda t}\) hasta una escala \(\Delta\). Calcule FTLE en tres
ventanas y muestre su distribución, no sólo la media.

Fije una región de escape, estime la supervivencia
de un conjunto de semillas y compare una ley exponencial con una cola no
exponencial. Trate las trayectorias que no escaparon como datos censurados.

Repita el diseño para Caputo en el espacio de historias:
defina las historias muestreadas, conserve el terminal inferior y separe
continuación de reinicio. Indique qué parte vive sólo en la proyección.
\end{bloqueestudio}
\fi}

\newcommand{\HAFOCardCases}{\ifHAFOPedagogy
\begin{bloqueestudio}{Lecturas}
Artículos primarios de los casos: Chua
\cite{LeonovKuznetsovVagaitsev2011,KuznetsovEtAl2017,
KuznetsovEtAl2023ChuaExperiments}; sistemas
fraccionarios \cite{Danca2017,Danca2017HiddenFO,Wu2023ArctanChua}; candidatos
no Chua \cite{munoz2018,wang2020}. La comparación de métodos se desarrolla en
\cite{DudkowskiEtAl2016,GuanXie2025}.
\end{bloqueestudio}
\begin{bloqueestudio}{Ejercicios}
Recalcule equilibrios, jacobiano y clasificación local de un
caso entero y uno fraccionario.

Para cada caso complete: campo; espacio de fases;
equilibrios; estabilidad; semilla; integración; destino; caos; cuenca;
conclusión matemática. Compare en paralelo un caso entero y uno fraccionario:
qué objetos permanecen en el espacio de estados y cuáles pasan al espacio de
historias.
\end{bloqueestudio}
\fi}

\newcommand{\HAFOCardCandidates}{\ifHAFOPedagogy
\begin{bloqueestudio}{Lecturas}
\emph{Systems with Hidden Attractors}
\cite{PhamVolosKapitaniak2017Systems}: pp.~21--63 y 79--102; catálogo
\cite{PhamEtAl2018HiddenAttractors}; capítulo fraccionario, pp.~199--238,
en \cite{WangKuznetsovChen2021Multistability}.
\end{bloqueestudio}
\begin{bloqueestudio}{Ejercicios}
Seleccione un candidato y recupere su artículo original.
Registre ecuaciones, parámetros, orden, inicialización y equilibrios.

Calcule jacobiano y estabilidad local; use el criterio
entero o de Matignon según corresponda. Después diseñe una semilla y una prueba
de cuenca que separen atracción, caos y ocultedad.
\end{bloqueestudio}
\fi}

\newcommand{\HAFOCardFuture}{\ifHAFOPedagogy
\begin{bloqueestudio}{Lecturas complementarias}
\begin{itemize}[nosep]
  \item \emph{Topología y geometría diferencial}: conexiones, pp.~135--142;
  geodésicas, pp.~143--149; torsión y curvatura, pp.~151--165; métrica,
  pp.~165--186.
  \item \emph{Teoría geométrica de ecuaciones diferenciales II}:
  formas normales y bifurcaciones, pp.~185--244.
  \item Guckenheimer--Holmes \cite{GuckenheimerHolmes1983}: PDF aprox.~6--39,
  65--119 y 120--184.
  \item Jost \cite{Jost2005Dynamical}: conjuntos invariantes aislados,
  pares índice y Conley, pp.~70--97; entropía y dinámica simbólica,
  pp.~99--110. Los pares índice permiten describir la persistencia de conjuntos invariantes aislados.
  \item \emph{Superfractals}, de Barnsley: PDF~19--99, 201--395;
  Lorenz, PDF~9--30 y 116--207; Mandelbrot, PDF~33, 277 y desde 367.
  \item Diethelm \cite{Diethelm2010AnalysisFDE}: multi-término,
  pp.~167--186. Oldham--Spanier \cite{OldhamSpanier1974Fractional}:
  pp.~120--132, 148--164, 192 y 207--218.
\end{itemize}
\end{bloqueestudio}
\begin{bloqueestudio}{Ejercicios}
Elija un caso entero y uno Caputo. Formule para ambos un problema abierto sobre
variedades, retorno o frontera de cuenca; escriba el objeto geométrico, el
espacio de fases correspondiente, la prueba numérica y la hipótesis que aún
faltaría demostrar.
\end{bloqueestudio}
\fi}
